\documentclass[10pt,a4paper]{amsart}
\usepackage{amsmath,amssymb,amsthm,graphicx}
\usepackage{verbatim}
\newenvironment{reference}{\comment}{\endcomment}
%\newenvironment{reference}{}{}
\newenvironment{slogan}{\comment}{\endcomment}
\newenvironment{history}{\comment}{\endcomment}

% For commutative diagrams we use Xy-pic
\usepackage[all]{xy}

% We use 2cell for 2-commutative diagrams.
\xyoption{2cell}
% Theorem environments.
%
\theoremstyle{plain}
\newtheorem{theorem}[subsection]{Theorem}
\newtheorem{proposition}[subsection]{Proposition}
\newtheorem{lemma}[subsection]{Lemma}

\theoremstyle{definition}
\newtheorem{definition}[subsection]{Definition}
\newtheorem{example}[subsection]{Example}
\newtheorem{exercise}[subsection]{Exercise}
\newtheorem{situation}[subsection]{Situation}

\theoremstyle{remark}
\newtheorem{remark}[subsection]{Remark}
\newtheorem{remarks}[subsection]{Remarks}

\numberwithin{equation}{subsection}

% Macros
%
\def\lim{\mathop{\mathrm{lim}}\nolimits}
\def\colim{\mathop{\mathrm{colim}}\nolimits}
\def\Spec{\mathop{\mathrm{Spec}}}
\def\Hom{\mathop{\mathrm{Hom}}\nolimits}
\def\Ext{\mathop{\mathrm{Ext}}\nolimits}
\def\SheafHom{\mathop{\mathcal{H}\!\mathit{om}}\nolimits}
\def\SheafExt{\mathop{\mathcal{E}\!\mathit{xt}}\nolimits}
\def\Sch{\mathit{Sch}}
\def\Mor{\mathop{\mathrm{Mor}}\nolimits}
\def\Ob{\mathop{\mathrm{Ob}}\nolimits}
\def\Sh{\mathop{\mathit{Sh}}\nolimits}
\def\NL{\mathop{N\!L}\nolimits}
\def\CH{\mathop{\mathrm{CH}}\nolimits}
\def\proetale{{pro\text{-}\acute{e}tale}}
\def\etale{{\acute{e}tale}}
\def\QCoh{\mathit{QCoh}}
\def\Ker{\mathop{\mathrm{Ker}}}
\def\Im{\mathop{\mathrm{Im}}}
\def\Coker{\mathop{\mathrm{Coker}}}
\def\Coim{\mathop{\mathrm{Coim}}}

% Boxtimes
%
\DeclareMathSymbol{\boxtimes}{\mathbin}{AMSa}{"02}

%
% Macros for moduli stacks/spaces
%
\def\QCohstack{\mathcal{QC}\!\mathit{oh}}
\def\Cohstack{\mathcal{C}\!\mathit{oh}}
\def\Spacesstack{\mathcal{S}\!\mathit{paces}}
\def\Quotfunctor{\mathrm{Quot}}
\def\Hilbfunctor{\mathrm{Hilb}}
\def\Curvesstack{\mathcal{C}\!\mathit{urves}}
\def\Polarizedstack{\mathcal{P}\!\mathit{olarized}}
\def\Complexesstack{\mathcal{C}\!\mathit{omplexes}}
% \Pic is the operator that assigns to X its picard group, usage \Pic(X)
% \Picardstack_{X/B} denotes the Picard stack of X over B
% \Picardfunctor_{X/B} denotes the Picard functor of X over B
\def\Pic{\mathop{\mathrm{Pic}}\nolimits}
\def\Picardstack{\mathcal{P}\!\mathit{ic}}
\def\Picardfunctor{\mathrm{Pic}}
\def\Deformationcategory{\mathcal{D}\!\mathit{ef}}
\usepackage{cleveref}
\crefname{originallemma}{original lemma}{original lemmas}
\Crefname{originallemma}{Original lemma}{Original lemmas}
\makeatletter
\newlabel{lemma-functor}{{0GNS}{source}{}{sourcepointer-1}{}}
\newlabel{lemma-functor@cref}{{[originallemma][0][]0GNS}{[1][0][]source}{}{}{}}
\makeatother
\crefname{originalremark}{original remark}{original remarks}
\Crefname{originalremark}{Original remark}{Original remarks}
\makeatletter
\newlabel{remark-composition}{{0GNW}{source}{}{sourcepointer-2}{}}
\newlabel{remark-composition@cref}{{[originalremark][0][]0GNW}{[1][0][]source}{}{}{}}
\makeatother
\crefname{originallemma}{original lemma}{original lemmas}
\Crefname{originallemma}{Original lemma}{Original lemmas}
\makeatletter
\newlabel{lemma-functor-modules}{{0GNV}{source}{}{sourcepointer-3}{}}
\newlabel{lemma-functor-modules@cref}{{[originallemma][0][]0GNV}{[1][0][]source}{}{}{}}
\makeatother
\crefname{originalreference}{original reference}{original references}
\Crefname{originalreference}{Original reference}{Original references}
\makeatletter
\newlabel{more-algebra-lemma-invertible}{{0B8I}{source}{}{sourcepointer-4}{}}
\newlabel{more-algebra-lemma-invertible@cref}{{[originalreference][0][]0B8I}{[1][0][]source}{}{}{}}
\makeatother
\crefname{originalreference}{original reference}{original references}
\Crefname{originalreference}{Original reference}{Original references}
\makeatletter
\newlabel{algebra-lemma-finite-projective}{{00NX}{source}{}{sourcepointer-5}{}}
\newlabel{algebra-lemma-finite-projective@cref}{{[originalreference][0][]00NX}{[1][0][]source}{}{}{}}
\makeatother
\title{An equivalence of module categories is tensoring with an invertible module over a ring isomorphism A → B}
\author{The Stacks Project Authors}
\date{stacks-project:functors.tex:local-snapshot-20261012; GNU FDL 1.2 or later}
\begin{document}
\maketitle
\noindent\textit{Editorial note.} All delivered mathematical text is the author's own, copied verbatim from the pinned Stacks Project chapter file functors.tex; no statement or proof step was rewritten, reordered or removed.\par\medskip
\noindent\textit{Editorial note.} Span recomputation by label pairing (LANE\_BRIEF 10.2): in this chapter the author's proof environment follows the statement environment as a standalone block; the spec covers the complete environment pair (statement span, proof span) as recorded in the pool candidate.\par\medskip
\noindent\textit{Editorial note.} Occupation scan: combined claim\_key + statement\_anchor + generic-line scan over all registration files (active and withdrawn) found 0 hits for functors.tex (fu-occupied.json, 2026-10-14); consistent with the duty manager's six-check dispatch.\par\medskip
\ifdefined\hypertarget\hypertarget{sourcepointer-1}{}\fi
\noindent\textit{Source reference.} Original lemma 0GNS: sources/stacks-chow/tags.txt, source lines 20802--20802. Cross-chapter standard prerequisite cited by the author inside the counted span; it is used as a pointer and its own proof is not part of this item.\par\medskip
\ifdefined\hypertarget\hypertarget{sourcepointer-2}{}\fi
\noindent\textit{Source reference.} Original remark 0GNW: sources/stacks-chow/tags.txt, source lines 20806--20806. Cross-chapter standard prerequisite cited by the author inside the counted span; it is used as a pointer and its own proof is not part of this item.\par\medskip
\ifdefined\hypertarget\hypertarget{sourcepointer-3}{}\fi
\noindent\textit{Source reference.} Original lemma 0GNV: sources/stacks-chow/tags.txt, source lines 20805--20805. Cross-chapter standard prerequisite cited by the author inside the counted span; it is used as a pointer and its own proof is not part of this item.\par\medskip
\ifdefined\hypertarget\hypertarget{sourcepointer-4}{}\fi
\noindent\textit{Source reference.} Original reference 0B8I: sources/stacks-chow/tags.txt, source lines 14141--14141. Cross-chapter standard prerequisite cited by the author inside the counted span; it is used as a pointer and its own proof is not part of this item.\par\medskip
\ifdefined\hypertarget\hypertarget{sourcepointer-5}{}\fi
\noindent\textit{Source reference.} Original reference 00NX: sources/stacks-chow/tags.txt, source lines 875--875. Cross-chapter standard prerequisite cited by the author inside the counted span; it is used as a pointer and its own proof is not part of this item.\par\medskip
\begin{lemma}
\label{lemma-functor-equivalence}
Let $A$ and $B$ be rings. If
$$
F : \text{Mod}_A \longrightarrow \text{Mod}_B
$$
is an equivalence of categories, then there exists an isomorphism
$A \to B$ of rings and an invertible $B$-module $L$ such that
$F$ is isomorphic to the functor $M \mapsto (M \otimes_A B) \otimes_B L$.
\end{lemma}

\begin{proof}
Since an equivalence commutes with all colimits, we see that
Lemmas \ref{lemma-functor} applies. Let $K$ be the
$A \otimes_\mathbf{Z} B$-module such that $F$ is
isomorphic to the functor $M \mapsto M \otimes_A K$.
Let $K'$ be the $B \otimes_\mathbf{Z} A$-module such that
a quasi-inverse of $F$ is
isomorphic to the functor $N \mapsto N \otimes_B K'$.
By Remark \ref{remark-composition} and
Lemma \ref{lemma-functor-modules} we have an isomorphism
$$
\psi : K \otimes_B K' \longrightarrow A
$$
of $A \otimes_\mathbf{Z} A$-modules.
Similarly, we have an isomorphism
$$
\psi' : K' \otimes_A K \longrightarrow B
$$
of $B \otimes_\mathbf{Z} B$-modules. Choose an element
$\xi = \sum_{i = 1, \ldots, n} x_i \otimes y_i \in K \otimes_B K'$
such that $\psi(\xi) = 1$. Consider the isomorphisms
$$
K \xrightarrow{\psi^{-1} \otimes \text{id}_K}
K \otimes_B K' \otimes_A K \xrightarrow{\text{id}_K \otimes \psi'} K
$$
The composition is an isomorphism and given by
$$
k \longmapsto \sum x_i \psi'(y_i \otimes k)
$$
We conclude this automorphism factors as
$$
K \to B^{\oplus n} \to K
$$
as a map of $B$-modules. It follows that $K$ is finite
projective as a $B$-module.

\medskip\noindent
We claim that $K$ is invertible as a $B$-module. This is equivalent
to asking the rank of $K$ as a $B$-module to have the constant value $1$,
see More on Algebra, Lemma \ref{more-algebra-lemma-invertible} and
Algebra, Lemma \ref{algebra-lemma-finite-projective}.
If not, then there exists a maximal ideal $\mathfrak m \subset B$
such that either (a) $K \otimes_B B/\mathfrak m = 0$ or
(b) there is a surjection $K \to (B/\mathfrak m)^{\oplus 2}$ of
$B$-modules. Case (a) is absurd as $K' \otimes_A K \otimes_B N = N$
for all $B$-modules $N$. Case (b) would imply we get a surjection
$$
A = K \otimes_B K' \longrightarrow (B/\mathfrak m \otimes_B K')^{\oplus 2}
$$
of (right) $A$-modules. This is impossible as the target is an $A$-module
which needs at least two generators: $B/\mathfrak m \otimes_B K'$
is nonzero as the image of the nonzero module $B/\mathfrak m$ under
the quasi-inverse of $F$.

\medskip\noindent
Since $K$ is invertible as a $B$-module we see that $\Hom_B(K, K) = B$.
Since $K = F(A)$ the action of $A$ on $K$ defines a ring isomorphism
$A \to B$. The lemma follows.
\end{proof}

\end{document}
